\documentclass[11pt]{book}
\usepackage[a4paper,textwidth=150mm]{geometry}
\usepackage[T1]{fontenc}
\usepackage{newpxtext}
\usepackage{microtype}
% font expansion with protrusion: a line ending in a ligature (T1's
% closing quotes) varies its marginal kern by pdfTeX's char_pw of the
% ligature's lig_char, which is no char node: 0, not the character's
% left protrusion (line_break's "Calculate variations of marginal
% kerns"); the item's lines break as pdfTeX's
\begin{document}
\begin{itemize}
\item \textbf{Cost of an implementation.} The only expensive step is the Kirchhoff solve. The paper points to nearly-linear-time Laplacian solvers (Spielman and Teng), to the correspondence between electrical networks and random walks (which would allow a fully decentralised solver), to auction algorithms, and to Baum's ``knee-jerk'' algorithm as possible relatives.
\item The paper points to the correspondence between electrical networks and random walks (which would allow a fully decentralised solver), to auction algorithms, and to Baum's knee-jerk'' algorithm as relatives, fine'' and final'' ones.
\end{itemize}
\end{document}
